\section*{Chapter \ref{ch:g9:income-distribution} --- Who Earns What: Income Distribution}

\begin{solution}{exo:g9:income-distribution:1}
Around 1990, at about 3 dollars a day (the highest band is centred on $2.8$). Around 2019, at about 6
dollars a day (the band centred on $5.6$).
\end{solution}

\begin{solution}{exo:g9:income-distribution:2}
Sum $8 + 9 + 10 + 10 + 11 + 12 + 13 + 15 + 40 = 128$; mean $128 \div 9 = 14.2$ thousand dalens. Median: the
fifth value, 11 thousand dalens.
\end{solution}

\begin{solution}{exo:g9:income-distribution:3}
Total $30 + 50 + 70 + 100 + 250 = 500$. Shares: $6\,\%$, $10\,\%$, $14\,\%$, $20\,\%$ and $50\,\%$.
\end{solution}

\begin{solution}{exo:g9:income-distribution:4}
Sweden 2023: about $37\,\%$ (table value $37.2$). Brazil 2024: about $55\,\%$ (table value $55.4$).
\end{solution}

\begin{solution}{exo:g9:income-distribution:5}
Mean of all: $49\,840\,000 \div 2\,000 = 24\,920$ dalens. Poorest fifth: $0.111 \times 49\,840\,000 =
5\,532\,240$ dalens shared by 400 \omterm{def:g7:households:household}{households}, $5\,532\,240 \div 400 = 13\,831$ dalens. (The survey's exact
value is 13,869: the share $11.1\,\%$ is rounded.)
\end{solution}

\begin{solution}{exo:g9:income-distribution:6}
New mean $(128 + 200) \div 10 = 32.8$ thousand dalens. Sorted, the fifth and sixth values are 11 and 12, so
the median is $11.5$. The mean rose by $18.6$, the median by $0.5$: the mean changed much more.
\end{solution}

\begin{solution}{exo:g9:income-distribution:7}
$43.7 - 37.2 = 6.5$ percentage points.
\end{solution}

\begin{solution}{exo:g9:income-distribution:8}
Before: $60 \div 6 = 10$ times. After: $72 \div 12 = 6$ times.
\end{solution}

\begin{solution}{exo:g9:income-distribution:9}
A higher share below the line in 2019 than in 1990, measured the same way. The chart's method could hide it
if the countries left out in 1990 (16\,\% of the world's people) had been much richer than those in, or if
the countries left out in 2019 were much poorer; or if surveys measured \omterm{def:g7:households:income}{income} in one year and consumption
in the other.
\end{solution}

\begin{solution}{exo:g9:income-distribution:10}
The richest fifth falls from 47 to $42\,\%$, the poorest rises from 6 to $11\,\%$: $11$, $10$, $15$, $22$,
$42\,\%$.
\end{solution}

\begin{solution}{exo:g9:income-distribution:11}
\begin{center}\small
\begin{tabular}{lrr}
\toprule
 & around 1990 (or 1980--81) & around 2019 (or 2022--24) \\
\midrule
World, share below 3 dollars a day & $45\,\%$ & under $10\,\%$ \\
United States, richest fifth's share & 41.2 (1980) & 47.1 (2024) \\
China, richest fifth's share & 37.2 (1981) & 43.7 (2022) \\
\bottomrule
\end{tabular}
\end{center}
Both are true because most of the world's poorest people lived in countries, China first, whose \omterm{def:g7:households:income}{incomes} grew
fast: even when the richest fifth's share rises, the poorest can see their \omterm{def:g7:households:income}{incomes} grow a lot.
\end{solution}

\begin{solution}{exo:g9:income-distribution:12}
No. The shares say how \omterm{def:g7:households:income}{income} is split among the five groups each year, not who is in each group. Some
\omterm{def:g7:households:household}{households} may have moved up and others down. One would need to follow the same \omterm{def:g7:households:household}{households} over the years
(a survey that interviews the same \omterm{def:g7:households:household}{households} again).
\end{solution}

\begin{solution}{pb:g9:income-distribution:1}
\textbf{1.} $49\,840\,000 \div 2\,000 = 24\,920$ dalens.
\textbf{2.} Below: $22\,265 < 24\,920$.
\textbf{3.} A few \omterm{def:g7:households:household}{households} with high \omterm{def:g7:households:income}{incomes} pull the mean up; the median does not move with them.
\textbf{4.} $37.4 \div 8.1 = 4.6$ times.
\textbf{5.} $33.5 \div 11.1 = 3.0$ times.
\textbf{6.} The poorest, second and middle fifths have larger shares ($8.1 \to 11.1$, $14.0 \to 15.0$, $17.5
\to 17.9$); the fourth and richest smaller ones ($23.0 \to 22.5$, $37.4 \to 33.5$).
\textbf{7.} $26\,818 - 3\,391 + 1\,493 = 24\,920$ dalens.
\textbf{8.} $3 \times 2 = 6$ dalens a day.
\textbf{9.} $6 \times 365 = 2\,190$ dalens a year.
\textbf{10.} $2 \times 2\,190 = 4\,380$ dalens a year.
\textbf{11.} $13\,869 \div 4\,380 = 3.2$ times.
\textbf{12.} No: a mean says nothing about the lowest \omterm{def:g7:households:income}{incomes}, and \omterm{def:g7:households:household}{households} differ in size. One needs each
\omterm{def:g7:households:household}{household}'s \omterm{def:g7:households:income}{income} and number of people, to compare \omterm{def:g6:living-standards:income-per-person}{income per person} with $2\,190$ dalens.
\textbf{13.} $37.4 - 33.5 = 3.9$ percentage points.
\end{solution}
